\begin{problem}[Seating around a circle]
How many different ways can $9$ people be seated in a circle if $2$ people have already chosen their seats?
\end{problem}

\begin{solution}
Once the two chosen seats are fixed, those two people no longer move. We only need to arrange the remaining $7$ people in the remaining $7$ seats.
\[
7! = 5040
\]
Therefore, the number of possible seatings is $\boxed{5040}$.
\end{solution}

\begin{takeaways}
\item Circular arrangements usually reduce by one factor because rotations are equivalent.
\item If some seats are already fixed, the circular symmetry has already been broken.
\end{takeaways}

\begin{problem}[Selections from URGOATED]
Five letters are chosen from the letters of the word \texttt{URGOATED}, without ordering them. How many selections are there if the choice:
\begin{enumerate}[label=\alph*)]
\item has no restriction?
\item must contain all four vowels?
\item must contain the letter \texttt{R}?
\item must contain the letter \texttt{U}, but not \texttt{D}?
\end{enumerate}
\end{problem}

\begin{solution}
All $8$ letters in \texttt{URGOATED} are distinct: $U, R, G, O, A, T, E, D$.

\textbf{a)} No restriction:
\[
\binom{8}{5} = 56
\]

\textbf{b)} Must contain all four vowels $U, O, A, E$:
We must choose $1$ more letter from the remaining $4$ consonants $R, G, T, D$.
\[
\binom{4}{1} = 4
\]

\textbf{c)} Must contain \texttt{R}:
Fix \texttt{R}, then choose $4$ more letters from the remaining $7$ letters.
\[
\binom{7}{4} = 35
\]

\textbf{d)} Must contain \texttt{U}, but not \texttt{D}:
Fix \texttt{U}, exclude \texttt{D}, then choose $4$ more letters from $R, G, O, A, T, E$.
\[
\binom{6}{4} = 15
\]

So the answers are:
\[
\boxed{56,\ 4,\ 35,\ 15}
\]
\end{solution}

\begin{takeaways}
\item Restriction questions are often easiest if you first force required letters in, or remove forbidden letters.
\item After handling the restriction, finish with a simple combination count.
\end{takeaways}

\begin{problem}[Forming a Dota 2 e-sport team]
\begin{enumerate}[label=\alph*)]
\item How many different ways can a Dota 2 e-sport team of $5$ students be chosen from a class of $30$ students?
\item The team must have $2$ girls and $3$ boys. If the class contains $12$ girls and $18$ boys, how many teams are possible?
\item Suppose Vu and Zachary are both students in the class. Find the probability that both Vu and Zachary are chosen on the team.
\item Find the probability that neither Vu nor Zachary is chosen.
\end{enumerate}
\end{problem}

\begin{solution}
\textbf{a)} Choose any $5$ students from $30$:
\[
\binom{30}{5} = 142506
\]

\textbf{b)} Choose $2$ girls from $12$ and $3$ boys from $18$:
\[
\binom{12}{2}\binom{18}{3} = 66 \times 816 = 53856
\]

\textbf{c)} If Vu and Zachary must both be chosen, then choose the remaining $3$ team members from the other $28$ students:
\[
\frac{\binom{28}{3}}{\binom{30}{5}} = \frac{3276}{142506} = \frac{2}{87}
\]
So the probability is $\boxed{\frac{2}{87}}$.

\textbf{d)} If neither Vu nor Zachary is chosen, then all $5$ team members must come from the other $28$ students:
\[
\frac{\binom{28}{5}}{\binom{30}{5}} = \frac{98280}{142506} = \frac{460}{667}
\]
So the probability is $\boxed{\frac{460}{667}}$.
\end{solution}

\begin{takeaways}
\item For probabilities in equally likely team selections, count favourable teams over total teams.
\item ``Both chosen'' means fix those people first; ``neither chosen'' means remove them from the pool.
\end{takeaways}

\begin{problem}[Triangles from points on the sides of a triangle]
The diagram shows triangle $ABC$ with points chosen on each of the sides. On side $AB$, $3$ points are chosen. On side $AC$, $4$ points are chosen. On side $BC$, $5$ points are chosen.

\begin{center}
\begin{tikzpicture}[scale=1.2]
    \coordinate (A) at (1,3);
    \coordinate (B) at (0,0);
    \coordinate (C) at (6,0.8);

    \draw (A) -- (B) -- (C) -- cycle;

    \node[above left] at (A) {$A$};
    \node[below left] at (B) {$B$};
    \node[right] at (C) {$C$};

    \foreach \i in {1,2,3}
        \filldraw [black] ($(A)!\i/4!(B)$) circle (1.5pt);

    \foreach \i in {1,2,3,4}
        \filldraw [black] ($(A)!\i/5!(C)$) circle (1.5pt);

    \foreach \i in {1,2,3,4,5}
        \filldraw [black] ($(B)!\i/6!(C)$) circle (1.5pt);
\end{tikzpicture}
\end{center}

How many triangles can be formed using the chosen points as vertices?
\end{problem}

\begin{solution}
There are
\[
3+4+5=12
\]
chosen points in total.

If no three points were collinear, the number of triangles would be
\[
\binom{12}{3}=220.
\]

However, any $3$ points chosen on the same side do not form a triangle.

On side $AB$, there are
\[
\binom{3}{3}=1
\]
collinear triple.

On side $AC$, there are
\[
\binom{4}{3}=4
\]
collinear triples.

On side $BC$, there are
\[
\binom{5}{3}=10
\]
collinear triples.

So the number of valid triangles is
\[
\binom{12}{3}-\binom{3}{3}-\binom{4}{3}-\binom{5}{3}
=220-1-4-10
=205.
\]

Hence the number of triangles is $\boxed{205}$.
\end{solution}

\begin{takeaways}
\item Start by counting all possible triples of points.
\item Then subtract the triples that are collinear, since they do not form triangles.
\item Geometric counting problems often reduce to combinations plus careful exclusion.
\end{takeaways}

\begin{problem}[Age-limit check across teams]
A sports association manages $13$ junior teams. It decides to check the age of all players. Any team that has more than $3$ players above the age limit will be penalised.

A total of $41$ players are found to be above the age limit.

Will any team be penalised? Justify your answer.
\end{problem}

\begin{solution}
Yes, at least one team must be penalised.

Suppose no team were penalised. Then each of the $13$ teams would have at most $3$ players above the age limit.

So the greatest possible total number of above-limit players would be
\[
13 \times 3 = 39.
\]

But the actual total is $41$, and
\[
41 > 39.
\]

This is impossible if every team has at most $3$ above-limit players.

Therefore, at least one team must have more than $3$ players above the age limit, so at least one team will be penalised.
\end{solution}

\begin{takeaways}
\item This is a pigeonhole-principle style argument: if every group had the maximum allowed amount, the total would still be too small.
\item A good way to prove that something \emph{must} happen is to assume it does not, then show the numbers do not work.
\item In combinatorics, ``justify your answer'' often means giving a short contradiction argument like this one.
\end{takeaways}

\begin{problem}[Expected score with replacement]
A game consists of randomly selecting $4$ balls from a bag. After each ball is selected it is replaced in the bag. The bag contains $3$ red balls and $7$ green balls.

For each red ball selected, $10$ points are earned and for each green ball selected, $5$ points are deducted. For instance, if a player picks $3$ red balls and $1$ green ball, the score will be
\[
3\times 10 - 1\times 5 = 25
\]
points.

What is the expected score in the game?
\end{problem}

\begin{solution}
For one draw,
\[
P(\text{red})=\frac{3}{10}, \qquad P(\text{green})=\frac{7}{10}.
\]

So the expected score from one draw is
\begin{align*}
E(\text{score per draw})
&= 10\left(\frac{3}{10}\right) + (-5)\left(\frac{7}{10}\right) \\
&= 3 - \frac{35}{10} \\
&= 3 - 3.5 \\
&= -0.5.
\end{align*}

Since there are $4$ draws with replacement, the expected total score is
\[
4\times (-0.5) = -2.
\]

Therefore, the expected score in the game is $\boxed{-2}$ points.
\end{solution}

\begin{takeaways}
\item With replacement, each draw has the same probabilities.
\item The expected value of a total is the sum of the expected values of the individual draws.
\item A negative expected value means that, on average, the game favours the house rather than the player.
\end{takeaways}

\begin{problem}[Minimum votes needed to win]
The members of a club voted for a new president. There were $15$ candidates for the position of president and $3543$ members voted. Each member voted for one candidate only.

One candidate received more votes than anyone else and so became the new president.

What is the smallest number of votes the new president could have received?
\end{problem}

\begin{solution}
To make the winner's vote total as small as possible, we want the other $14$ candidates to receive as many votes as possible, while still having fewer votes than the winner.

Suppose the winner received $x$ votes. Then each of the other $14$ candidates can receive at most $x-1$ votes.

So the greatest possible total number of votes is
\[
x + 14(x-1).
\]
This must be at least $3543$:
\begin{align*}
x + 14(x-1) &\ge 3543 \\
15x - 14 &\ge 3543 \\
15x &\ge 3557 \\
x &\ge \frac{3557}{15} = 237.13\ldots
\end{align*}

So the smallest possible integer value of $x$ is
\[
\boxed{238}.
\]

This value is achievable, since if the winner gets $238$ votes then the other $14$ candidates can together receive up to
\[
14\times 237 = 3318
\]
votes, and
\[
3318 + 238 = 3556 > 3543,
\]
so the $3543$ votes can be distributed with the winner still strictly ahead.
\end{solution}

\begin{takeaways}
\item To minimize a maximum, try making all competitors as large as possible without overtaking it.
\item This is a classic ``worst-case'' or extremal counting argument.
\item After finding the bound, always check that it is actually achievable.
\end{takeaways}

\begin{problem}[Selecting finalists and awarding places]
Out of $10$ contestants, six are to be selected for the final round of a competition. Four of those six will be placed $1$st, $2$nd, $3$rd and $4$th.

In how many ways can this process be carried out?
\end{problem}

\begin{solution}
First choose the $6$ finalists from the $10$ contestants:
\[
\binom{10}{6} = 210.
\]

Now choose and order the $4$ placed contestants from those $6$ finalists:
\[
{}^6P_4 = 6\times 5\times 4\times 3 = 360.
\]

So the total number of ways is
\[
\binom{10}{6}\times {}^6P_4 = 210\times 360 = 75600.
\]

Hence the process can be carried out in $\boxed{75600}$ ways.
\end{solution}

\begin{takeaways}
\item Break the process into stages: first select, then arrange.
\item Use combinations for choosing the finalists because order does not matter there.
\item Use permutations for the placings because $1$st, $2$nd, $3$rd, and $4$th are ordered positions.
\end{takeaways}

\begin{problem}[Using the pigeonhole principle with topic choices]
To complete a course, a student must choose and pass exactly three topics. There are eight topics from which to choose.

Last year $400$ students completed the course.

Explain, using the pigeonhole principle, why at least eight students passed exactly the same three topics.
\end{problem}

\begin{solution}
The number of different ways to choose exactly $3$ topics from $8$ is
\[
\binom{8}{3} = 56.
\]

So there are only $56$ possible topic combinations.

Now distribute the $400$ students among these $56$ combinations. If at most $7$ students had any one particular combination, then the total number of students could be at most
\[
56\times 7 = 392.
\]

But in fact there are $400$ students, and
\[
400 > 392.
\]

Therefore, by the pigeonhole principle, at least one of the $56$ topic combinations must have been chosen by at least $8$ students.

So at least eight students passed exactly the same three topics.
\end{solution}

\begin{takeaways}
\item Start by counting the number of possible categories, or ``pigeonholes''.
\item If the total number of objects exceeds the maximum possible without overlap, then one category must contain more than that maximum.
\item In these questions, the pigeonhole principle often appears as: ``If every group had at most $k$, the total would be too small.'' 
\end{takeaways}
